\subsection{Adaptive Quizzing}
\label{sec:impl_quizzing}

Section~\ref{sec:impl_ai_pipelines} covers how quiz questions are generated. This section covers what happens afterwards: review, publication, the attempt itself, grading, and release. The implementation is in \path{abridgeai/features/quizzes/}, whose service layer is deliberately one module per concern --- \path{publish_gate.py}, \path{grader.py}, \path{gradebook.py}, \path{regrade.py}, \path{overrides.py}, \path{integrity.py}, \path{review_visibility.py} --- because these rules change independently of one another. Figure~\ref{fig:quiz_attempt_lifecycle} in Appendix~\ref{appendix:algorithm} shows the attempt itself end to end, where these concerns --- integrity monitoring, spaced repetition and recovery from an abandoned attempt --- meet on the same rows.

\paragraph{Human review before publication}
Every question carries a \texttt{review\_status} of \texttt{pending}, \texttt{approved}, \texttt{edited}, or \texttt{rejected}. AI-drafted questions enter as \texttt{pending} and only approved questions are ever delivered in an attempt. \path{publish_gate.py} refuses publication while any question is unreviewed, and separately refuses it while any question lacks an expected answer time --- the latter because the spaced-repetition scheduler grades recall partly on elapsed time against expectation, so a question with no expectation cannot be scheduled and must not reach a learner.

\paragraph{Per-student rule resolution}
A quiz's dates, time limit, and retake rules can be adjusted for an individual or a group. \path{override_policy.py} resolves the effective rules for one learner in a fixed precedence: the quiz's own settings, then the most generous group adjustment that applies to them, then their personal adjustment. Group adjustments resolve to the most generous rather than the last-written, because an accommodation must never be narrowed by a second one the learner also qualifies for.

\paragraph{The attempt}
An attempt moves through \texttt{in\_progress} to one of \texttt{submitted}, \texttt{graded}, \texttt{abandoned}, or \texttt{expired}. Question and option order are shuffled per attempt by \path{shuffle.py} from a per-attempt seed, so the order is stable across a reload but differs between learners. An attempt still open when its time limit or the quiz's close date passes is resolved automatically --- submitted, granted a grace period, or closed ungraded --- as the teacher configured, so an attempt abandoned by a closed laptop does not wait on the learner's return.

\paragraph{Grading and the grade of record}
\path{grader.py} marks objectively-answerable questions on submission and sets aside anything it cannot mark with confidence for the queue in \path{manual_grading.py}. While any answer in an attempt remains unmarked the learner is shown no score, since a partial score presented as final would mislead. \path{gradebook.py} then reduces a learner's attempts to one grade of record by the quiz's declared method --- highest, first, last, or average --- and stores the method and the attempts counted alongside the figure, so a grade can be explained later rather than merely asserted.

\paragraph{Regrading}
\path{regrade.py} retains an internal dry-run and commit foundation for re-evaluating affected attempts, but the capability is not exposed in the current teacher interface. Published questions are frozen and there is no controlled answer-key-correction workflow; exposing regrade without that entry point would produce an action that normally reports no changes, while applying a correction without reconciling spaced-repetition state could leave grades and ease factors inconsistent. The complete correction-to-regrade workflow is therefore deferred to future work.

\paragraph{Integrity signals}
\path{integrity.py} records window blur, tab switches, full-screen exits, and connection loss while an attempt is running. Recording never blocks or fails a submission, signals arriving after submission are discarded, and the weights and flag threshold are fixed at the moment the attempt starts --- so editing a quiz mid-cohort cannot change how an attempt already under way is judged. The records are evidence for a teacher to review, not a verdict, and learners never see them.

\paragraph{One client per attempt}
\path{session_guard.py} binds a live attempt to the authentication session that owns it, through the Redis key \path|quiz_attempt_session:{attempt_id}| held for ninety seconds. Claim, heartbeat, takeover and release are compare-and-set Lua scripts, so two tabs resuming the same attempt cannot both win the race; the frontend additionally holds a Web Locks key per attempt and coordinates handover over a \path{BroadcastChannel}, falling back to a channel leader handshake where Web Locks is unavailable. A learner moving to another device confirms an explicit takeover.

Two decisions here run against the grain of the rest of the backend and are deliberate. First, answer saves, integrity events and submission \textit{fail closed} when Redis is unreachable, where every other Redis path degrades to PostgreSQL --- admitting two concurrent clients to a graded attempt is worse than refusing the write. Second, the key is ephemeral rather than a durable lease: an abandoned client releases the attempt by falling silent. The server keeps the authoritative timer throughout, so closing a tab neither pauses nor forfeits an attempt; an accidental close leaves it resumable, while a deliberate exit flushes the current answer, releases the key, and stops the camera.

The boundary is worth stating plainly. A copied access token reusing the same authentication session id is indistinguishable from the original client, Redis loss or an administrative flush drops ownership, and a modified client can ignore the frontend tab coordination entirely. These are accepted limits of a session-scoped guard; closing them needs a persistent assessment lease shared by quizzes and interviews, with opaque per-client credentials, generation counters, audit history, invalidation on a Redis reset, and server-side enforcement reaching into the interview transport. That is left to future work.

\paragraph{Required camera}
\texttt{Quiz.require\_camera} defaults to false. Where a teacher enables it, the learner must grant \path{getUserMedia} video access before starting or resuming, and the workspace blocks whenever the video track is missing, ended, or muted --- after a four-second grace period, so a momentary device glitch does not interrupt the attempt. The timer and the ownership heartbeat both continue during that block, because the requirement governs whether the learner may keep working, not whether the clock runs.

As implemented, the check is a presence test and nothing more: the requirement is satisfied by an active video track, and the stream is held only as a local \path{MediaStream} in the browser --- no capture, upload, or analysis path exists in the current code. The flag is copied into the attempt's \texttt{integrity\_policy\_snapshot} at start alongside the signal weights, so a quiz edited later cannot retroactively impose --- or lift --- a camera requirement on an attempt already under way. Doing anything with the video itself, and the privacy and consent obligations that would follow, is left to future work.
